\documentclass[11pt,a4paper]{article}

% ------------------------------------------------------------------
%  Ultime Team Manager - Documento de planificacion
%  Compilar con:  tectonic planificacion.tex
% ------------------------------------------------------------------

\usepackage[T1]{fontenc}
\usepackage[utf8]{inputenc}
\usepackage[spanish,es-noshorthands]{babel}
\usepackage[a4paper,margin=2.2cm]{geometry}
\usepackage{xcolor}
\usepackage{colortbl}
\usepackage{booktabs}
\usepackage{array}
\usepackage{longtable}
\usepackage{tabularx}
\usepackage{enumitem}
\usepackage{titlesec}
\usepackage{fancyhdr}
\usepackage{pgfgantt}
\usepackage{hyperref}

% ---- Paleta ------------------------------------------------------
\definecolor{brand}{HTML}{0B7A4B}      % verde cesped
\definecolor{brandDark}{HTML}{064D2F}
\definecolor{accent}{HTML}{1565C0}     % azul
\definecolor{gold}{HTML}{C9A227}       % monedas
\definecolor{lightgray}{HTML}{F1F3F4}
\definecolor{midgray}{HTML}{6B7280}

\hypersetup{
  colorlinks=true,
  linkcolor=brandDark,
  urlcolor=accent,
  pdftitle={Ultime Team Manager - Planificacion},
  pdfauthor={Sebastian Ramirez, Joaquin Parraud}
}

% ---- Encabezado / pie -------------------------------------------
\pagestyle{fancy}
\fancyhf{}
\renewcommand{\headrulewidth}{0.4pt}
\lhead{\small\textcolor{midgray}{Ultime Team Manager}}
\rhead{\small\textcolor{midgray}{Planificacion}}
\cfoot{\small\textcolor{midgray}{\thepage}}

% ---- Titulos de seccion -----------------------------------------
\titleformat{\section}
  {\normalfont\Large\bfseries\color{brandDark}}
  {\thesection.}{0.6em}{}
\titleformat{\subsection}
  {\normalfont\large\bfseries\color{accent}}
  {\thesubsection}{0.6em}{}

% ---- Columnas de tabla ------------------------------------------
\newcolumntype{L}[1]{>{\raggedright\arraybackslash}p{#1}}
\newcolumntype{C}[1]{>{\centering\arraybackslash}p{#1}}

% ---- Caja de portada --------------------------------------------
\newcommand{\coverbox}{%
  \begingroup
  \setlength{\fboxsep}{18pt}%
  \noindent\colorbox{brand}{%
    \begin{minipage}{\dimexpr\textwidth-2\fboxsep\relax}
      \color{white}
      {\Huge\bfseries Ultime Team Manager}\\[6pt]
      {\large Plan de proyecto -- Unidad de Flutter}\\[2pt]
      {\normalsize Aplicacion tipo FC Ultimate Team en Flutter / Dart}
    \end{minipage}%
  }%
  \endgroup
}

\begin{document}

% ================================================================
%  PORTADA
% ================================================================
\thispagestyle{empty}
\vspace*{1.2cm}
\coverbox

\vspace{1.4cm}
\noindent
\renewcommand{\arraystretch}{1.4}
\begin{tabularx}{\textwidth}{@{}l X@{}}
\textbf{Integrantes} & Sebastian Ramirez \quad\textbullet\quad Joaquin Parraud \\
\textbf{Repositorio} & \texttt{Ultime-Team-Manager} (GitHub) \\
\textbf{Stack} & Flutter, Dart, Riverpod, go\_router, Hive, Supabase (futuro) \\
\textbf{Fecha de inicio} & Lunes 13 de julio (Clase 2) \\
\textbf{Entrega / presentacion} & Semana del 20 al 25 de julio \\
\end{tabularx}

\vfill
\noindent\textcolor{midgray}{\rule{\textwidth}{0.4pt}}\\[4pt]
\noindent{\small\textcolor{midgray}%
{Documento vivo: se actualiza al cerrar cada clase. Genera el PDF con \texttt{tectonic planificacion.tex}.}}

\newpage
\tableofcontents
\newpage

% ================================================================
\section{Resumen y contexto}
% ================================================================
\textbf{Ultime Team Manager} es una aplicacion movil inspirada en \emph{FC Ultimate
Team}. El usuario inicia sesion, arma su equipo de once colocando cartas de
futbolistas en sus posiciones (formacion \textbf{4-3-3} por defecto), compra y vende
cartas en un mercado, simula partidos cuyo resultado depende de la valoracion media
de cada equipo y compite en una liga que reparte monedas segun el resultado.

El proyecto se construye sobre la app base de la Unidad~1 (un contador con
\texttt{setState}) y evoluciona hacia una arquitectura por capas. El desarrollo es
\textbf{offline-first} pero disenado para escalar a un backend online (Supabase) sin
reescribir la capa de presentacion.

\vspace{4pt}
\noindent\fcolorbox{gold}{lightgray}{%
  \begin{minipage}{\dimexpr\textwidth-2\fboxsep\relax}
  \textbf{Meta de la unidad:} llevar la idea a una app funcional en tres clases y
  presentarla la semana del 20 al 25 de julio, cumpliendo los requisitos minimos
  (varias pantallas + navegacion, una lista con CRUD, conexion a datos, UI cuidada y
  repositorio en GitHub con ramas por tarea).
  \end{minipage}}

% ================================================================
\section{Objetivos y alcance}
% ================================================================
\subsection*{Objetivo general}
Entregar una app Flutter funcional y presentable que implemente el ciclo completo:
\emph{login $\rightarrow$ armar plantilla $\rightarrow$ mercado $\rightarrow$ simular
partido $\rightarrow$ liga y monedas}.

\subsection*{Dentro del alcance}
\begin{itemize}[leftmargin=1.2em]
  \item Login con validacion de campos y sesion persistente (offline).
  \item Plantilla de 11 jugadores en formacion 4-3-3 sobre un campo.
  \item Cartas con valoracion y atributos; datos desde JSON semilla y API publica.
  \item Mercado con CRUD (comprar/vender) sobre la coleccion del usuario.
  \item Simulacion de partidos por valoracion media y liga con reparto de monedas.
  \item Estados de carga y de vacio en todas las listas.
\end{itemize}

\subsection*{Fuera del alcance (por ahora)}
\begin{itemize}[leftmargin=1.2em]
  \item Backend online real (queda la arquitectura lista para Supabase).
  \item Multijugador en tiempo real, notificaciones push y micropagos.
  \item Motor de simulacion minuto a minuto (se usa un modelo probabilistico simple).
\end{itemize}

% ================================================================
\section{Stack tecnologico}
% ================================================================
\renewcommand{\arraystretch}{1.3}
\begin{tabularx}{\textwidth}{@{}L{4.2cm} X@{}}
\toprule
\textbf{Area} & \textbf{Eleccion} \\
\midrule
Framework / lenguaje & Flutter, Dart (\texttt{sdk \^{}3.12.2}) \\
Gestion de estado    & Riverpod (\texttt{flutter\_riverpod}) \\
Navegacion           & \texttt{go\_router} \\
Persistencia local   & Hive (cartas, plantilla, liga) + \texttt{shared\_preferences} (sesion) \\
Backend online (futuro) & Supabase (Auth + Postgres) \\
Origen de cartas     & JSON semilla en \texttt{assets/data} + API publica de futbol (TheSportsDB) \\
Testing / Linting    & \texttt{flutter\_test} + \texttt{flutter\_lints} \\
Control de versiones & Git + GitHub (ramas por tarea, pull requests) \\
\bottomrule
\end{tabularx}

% ================================================================
\section{Roles y responsabilidades}
% ================================================================
Dos integrantes con areas principales; el trabajo se reparte por \emph{feature} y se
integra mediante pull requests para asegurar commits de ambos.

\renewcommand{\arraystretch}{1.3}
\begin{tabularx}{\textwidth}{@{}L{4cm} X@{}}
\toprule
\textbf{Integrante} & \textbf{Responsabilidades} \\
\midrule
\textbf{Sebastian Ramirez} & Arquitectura base y navegacion, login + validacion,
sesion, entidades del dominio, pantalla de plantilla (4-3-3), tema/UI. \\
\addlinespace
\textbf{Joaquin Parraud} & Mercado (CRUD), consumo de API y modelos de datos,
simulacion de partidos, liga y economia de monedas, estados de carga/vacio. \\
\bottomrule
\end{tabularx}

% ================================================================
\section{Hitos del proyecto}
% ================================================================
\renewcommand{\arraystretch}{1.35}
\begin{tabularx}{\textwidth}{@{}C{2.4cm} L{4.6cm} X@{}}
\toprule
\textbf{Fecha} & \textbf{Hito} & \textbf{Entregable} \\
\midrule
Lun 06-07 & \textbf{H0} -- Arranque & Grupo, idea definida y repositorio creado (hecho). \\
Lun 13-07 & \textbf{H1} -- Login con validacion & Login funcional con validacion, sesion y navegacion entre 3+ pantallas. \\
Dom 19-07 & \textbf{H2} -- Plantilla y datos & Modelos, plantilla 4-3-3, persistencia local y primeras cartas. \\
Lun 20-07 & \textbf{H3} -- API + tarjetas & Consumo de API mostrado en tarjetas dinamicas; mercado con CRUD (cierre de proyecto). \\
20--25 jul & \textbf{H4} -- Presentacion & Simulacion, liga, monedas, UI pulida, pruebas y presentacion grupal. \\
\bottomrule
\end{tabularx}

% ================================================================
\section{Cronograma (Gantt)}
% ================================================================
Del 13 al 25 de julio, organizado en dos sprints alineados con las clases.

\vspace{6pt}
\noindent\hspace*{-0.3cm}
\begin{ganttchart}[
    hgrid, vgrid,
    x unit=0.92cm,
    y unit chart=0.62cm,
    y unit title=0.7cm,
    title height=1,
    canvas/.append style={fill=lightgray, draw=midgray, line width=0.4pt},
    bar/.append style={fill=brand, draw=brandDark, rounded corners=2pt},
    bar height=0.55,
    bar label font=\small,
    group/.append style={fill=accent, draw=accent!60!black},
    group label font=\small\bfseries,
    milestone/.append style={fill=gold, draw=gold!60!black},
    milestone label font=\small\itshape,
    title/.append style={fill=brandDark},
    title label font=\small\bfseries\color{white},
    include title in canvas=false
  ]{1}{13}

  \gantttitle{Sprint 1 (13--19 jul)}{7}
  \gantttitle{Sprint 2 (20--25 jul)}{6} \\
  \gantttitlelist{13,14,15,16,17,18,19,20,21,22,23,24,25}{1} \\

  \ganttgroup{Sprint 1: Login + base}{1}{7} \\
  \ganttbar{Arquitectura y navegacion}{1}{2} \\
  \ganttbar{Login + validacion (H1)}{1}{2} \\
  \ganttmilestone{H1}{2} \\
  \ganttbar{Entidades y modelos}{3}{4} \\
  \ganttbar{Plantilla 4-3-3}{4}{6} \\
  \ganttbar{Persistencia local}{5}{7} \\
  \ganttmilestone{H2}{7} \\

  \ganttgroup{Sprint 2: Datos + juego}{8}{13} \\
  \ganttbar{API + tarjetas (H3)}{8}{9} \\
  \ganttbar{Mercado (CRUD)}{8}{10} \\
  \ganttmilestone{H3}{8} \\
  \ganttbar{Simulacion de partidos}{9}{11} \\
  \ganttbar{Liga y monedas}{10}{12} \\
  \ganttbar{UI, estados y pruebas}{11}{13} \\
  \ganttbar{Presentacion}{13}{13} \\
  \ganttmilestone{H4}{13}
\end{ganttchart}

\vspace{4pt}
\noindent{\small\textcolor{midgray}{Los numeros del eje corresponden a los dias de
julio. Los rombos dorados marcan los hitos H1--H4.}}

% ================================================================
\section{Desglose de tareas por sprint}
% ================================================================

\subsection{Sprint 1 -- Login y base (13--19 jul)}
\renewcommand{\arraystretch}{1.3}
\begin{longtable}{@{}C{1.2cm} L{7.6cm} L{3.4cm} C{2.2cm}@{}}
\toprule
\textbf{ID} & \textbf{Tarea} & \textbf{Responsable} & \textbf{RF} \\
\midrule
\endhead
T1.1 & Configurar arquitectura por capas y \texttt{go\_router} & Sebastian & --- \\
T1.2 & Renombrar proyecto y anadir dependencias (Riverpod, Hive) & Sebastian & --- \\
T1.3 & Pantalla de login con validacion de campos & Sebastian & RF1 \\
T1.4 & Sesion persistente con \texttt{shared\_preferences} & Sebastian & RF1 \\
T1.5 & Entidades del dominio (Player, Card, Team, Match, League) & Joaquin & RF3 \\
T1.6 & Semilla JSON de jugadores en \texttt{assets/data} & Joaquin & RF3 \\
T1.7 & Pantalla de plantilla con formacion 4-3-3 & Sebastian & RF2 \\
T1.8 & Persistencia local de la plantilla (Hive) & Joaquin & RF2 \\
\bottomrule
\end{longtable}

\subsection{Sprint 2 -- Datos y juego (20--25 jul)}
\renewcommand{\arraystretch}{1.3}
\begin{longtable}{@{}C{1.2cm} L{7.6cm} L{3.4cm} C{2.2cm}@{}}
\toprule
\textbf{ID} & \textbf{Tarea} & \textbf{Responsable} & \textbf{RF} \\
\midrule
\endhead
T2.1 & Consumir API publica de futbol y mapear a modelos & Joaquin & RF3 \\
T2.2 & Mostrar cartas dinamicas con estados de carga/vacio & Joaquin & RF3 \\
T2.3 & Mercado: comprar y vender (CRUD completo) & Joaquin & RF4 \\
T2.4 & Motor de simulacion por valoracion media & Sebastian & RF5 \\
T2.5 & Liga: calendario, clasificacion y reparto de monedas & Joaquin & RF6 \\
T2.6 & Dashboard de inicio (monedas y accesos) & Sebastian & --- \\
T2.7 & Pulido de UI, tema y widgets reutilizables & Sebastian & --- \\
T2.8 & Pruebas de widgets/logica y preparacion de la demo & Ambos & --- \\
\bottomrule
\end{longtable}

% ================================================================
\section{Estrategia de ramas Git}
% ================================================================
\begin{itemize}[leftmargin=1.2em]
  \item \texttt{main} -- rama estable y presentable.
  \item \texttt{develop} -- integracion continua del trabajo del equipo.
  \item \texttt{feature/<tarea>} -- una rama por tarea (p.\,ej. \texttt{feature/login-validacion}).
  \item Cada \emph{feature} se integra mediante \textbf{pull request} revisado por el otro integrante.
\end{itemize}
Objetivo: historial claro con \textbf{commits de ambos integrantes} y trazabilidad
tarea~$\rightarrow$~rama~$\rightarrow$~PR.

% ================================================================
\section{Riesgos y mitigaciones}
% ================================================================
\renewcommand{\arraystretch}{1.3}
\begin{tabularx}{\textwidth}{@{}L{5.4cm} C{1.8cm} X@{}}
\toprule
\textbf{Riesgo} & \textbf{Impacto} & \textbf{Mitigacion} \\
\midrule
Plazo muy ajustado (2 semanas) & Alto & Priorizar RF minimos; dejar simulacion/liga como incremento final. \\
API externa caida o con limites & Medio & Caer a la semilla JSON local (offline-first). \\
Conflictos de merge & Medio & Ramas por tarea y PRs pequenos y frecuentes. \\
Alcance excesivo (scope creep) & Medio & Congelar el alcance tras H3; extras solo si sobra tiempo. \\
Curva de Riverpod/Hive & Bajo & Repartir por area y apoyarse en ejemplos de clase. \\
\bottomrule
\end{tabularx}

% ================================================================
\section{Definicion de Hecho (DoD) y rubrica}
% ================================================================
Una tarea esta \textbf{Hecha} cuando: compila sin errores, pasa el \emph{linter},
tiene los estados de carga/vacio cuando aplica, y esta integrada en \texttt{main} via
PR revisado.

\vspace{4pt}
\noindent\textbf{Mapeo con los requisitos minimos de la unidad:}
\renewcommand{\arraystretch}{1.3}
\begin{tabularx}{\textwidth}{@{}L{6.2cm} X@{}}
\toprule
\textbf{Requisito de la unidad} & \textbf{Como se cumple} \\
\midrule
Varias pantallas + navegacion & Login, Home, Plantilla, Mercado, Liga y Partido (go\_router). \\
Una lista con CRUD & Mercado y plantilla: crear, ver, editar y eliminar cartas. \\
Conexion a datos & API publica de futbol y/o Supabase, con caché local. \\
UI cuidada con widgets & Layouts, cartas, campo de juego y estados de carga/vacio. \\
Repositorio en GitHub & Ramas por tarea, PRs y commits de ambos integrantes. \\
\bottomrule
\end{tabularx}

\vfill
\noindent\textcolor{midgray}{\rule{\textwidth}{0.4pt}}\\[2pt]
\noindent{\small\textcolor{midgray}{Ultime Team Manager -- Sebastian Ramirez y Joaquin Parraud -- Unidad de Flutter.}}

\end{document}
